
@@M0@@

Si @@M1@@ est un @@M2@@-anneau, on a donc pour tous @@M3@@ :
@@M4@@

Le @@M5@@-ième polynôme universel de la loi @@M6@@ sera noté @@M7@@~; c'est un polynôme de @@M8@@, dont les composantes isobares sont de poids variant de @@M9@@ à @@M10@@ par rapport à l'ensemble des variables @@M11@@, (1.15) et symétrique par rapport à ces deux groupes de variables. Pour @@M12@@, on a donc :
@@M13@@

Notons enfin que la loi @@M14@@ admet un élément unité, qui est @@M15@@, soit @@M16@@~; il sera parfois noté @@M17@@.

3.6. Les «~opérations @@M18@@~» sont définies par le diagramme commutatif :
@@M19@@
i.e. par le morphisme composé de foncteurs :
@@M20@@
et seront notées @@M21@@. Elles sont définies par leur valeur pour @@M22@@~; on a :
